% =============================================================
% referencias.tex — Bibliografía IEEE separada
%
% Referencias reales y verificables sobre:
%   - Criba de Eratóstenes y algoritmos de criba
%   - MPI y programación paralela
%   - Computación distribuida
%
% Compilar:
%   pdflatex referencias.tex
% =============================================================

\documentclass[12pt]{article}

% Nota: para hifenación en español instalar: sudo dnf install texlive-babel-spanish
\usepackage{babel}
\usepackage[utf8]{inputenc}
\usepackage[T1]{fontenc}
\usepackage{url}
\usepackage[margin=2.5cm]{geometry}

\title{Referencias Bibliográficas\\
       \large Criba de Eratóstenes Distribuida con MPI}
\author{}
\date{}

\begin{document}

\maketitle
\thispagestyle{empty}

% =============================================================
% Formato IEEE para bibliografías:
% [número]  Inicial(es). Apellido(s), "Título del artículo",
%           Nombre de la publicación, vol. X, no. Y,
%           pp. AAA--BBB, Mes Año.
% =============================================================

\begin{thebibliography}{10}

% ── Libros sobre MPI ─────────────────────────────────────────

\bibitem{gropp1999}
W.~Gropp, E.~Lusk, and A.~Skjellum,
\textit{Using MPI: Portable Parallel Programming with the Message Passing
Interface}, 2nd~ed.
Cambridge, MA, USA: MIT Press, 1999, ISBN 978-0-262-57134-5.

\bibitem{gropp1998advanced}
W.~Gropp, E.~Lusk, and R.~Thakur,
\textit{Using MPI-2: Advanced Features of the Message Passing Interface}.
Cambridge, MA, USA: MIT Press, 1999, ISBN 978-0-262-57133-8.

\bibitem{quinn2003}
M.~J. Quinn,
\textit{Parallel Programming in C with MPI and OpenMP}.
New York, NY, USA: McGraw-Hill Education, 2003, ISBN 978-0-072-82256-4.

\bibitem{pacheco2011}
P.~S. Pacheco,
\textit{An Introduction to Parallel Programming}.
Burlington, MA, USA: Morgan Kaufmann, 2011, ISBN 978-0-123-74260-5.

% ── Estándar MPI ─────────────────────────────────────────────

\bibitem{mpiforum2021}
Message Passing Interface Forum,
``MPI: A Message-Passing Interface Standard, Version 4.0,''
Technical Report, Jun.~2021.
[Online]. Available: \url{https://www.mpi-forum.org/docs/mpi-4.0/mpi40-report.pdf}

% ── Algoritmos numéricos y criba ─────────────────────────────

\bibitem{knuth1997}
D.~E. Knuth,
\textit{The Art of Computer Programming, Vol.~2: Seminumerical Algorithms},
3rd~ed.
Reading, MA, USA: Addison-Wesley, 1997, ISBN 978-0-201-89684-8.

\bibitem{bach1996}
E.~Bach and J.~Shallit,
\textit{Algorithmic Number Theory, Vol.~1: Efficient Algorithms}.
Cambridge, MA, USA: MIT Press, 1996, ISBN 978-0-262-02405-1.

\bibitem{bays1977}
C.~Bays and R.~H. Hudson,
``The segmented sieve of Eratosthenes and primes in arithmetic progressions
to $10^{12}$,''
\textit{BIT Numerical Mathematics}, vol.~17, no.~2, pp.~121--127, Jun.~1977.
DOI: \url{https://doi.org/10.1007/BF01932283}

\bibitem{pritchard1982}
P.~Pritchard,
``Explaining the wheel factorization,''
\textit{BIT Numerical Mathematics}, vol.~22, no.~4, pp.~442--454, Dec.~1982.
DOI: \url{https://doi.org/10.1007/BF01934411}

\bibitem{oliveira2014}
T.~Oliveira~e~Silva, S.~Herzog, and S.~Pardi,
``Empirical verification of the even Goldbach conjecture, and computation of
prime gaps, up to $4 \times 10^{18}$,''
\textit{Mathematics of Computation}, vol.~83, no.~288, pp.~2033--2060, 2014.
DOI: \url{https://doi.org/10.1090/S0025-5718-2013-02787-1}

\end{thebibliography}

\end{document}
